% Inserir na Discussão, após os resultados de generalização entre períodos.

A interpretação dos resultados é limitada pelo número de dez vasos e pela
associação entre dia de coleta, condição do solo e grupo experimental.
Houve irrigação entre 13 e 18 de maio de 2026, coincidindo com a mudança
da faixa operacional do índice de umidade do solo. A validação LOSO
separou vasos, mas manteve dias em comum entre treinamento e teste;
quando os períodos foram separados, o desempenho caiu. Não foram
documentadas para esta análise a randomização da posição dos vasos e da
ordem de medição, a aplicação aos controles do mesmo veículo sem J2,
nem a confirmação nematológica individual. Assim, não é possível
descartar efeitos de posição, sequência de aquisição, veículo ou deriva
instrumental como explicações parciais dos padrões observados. A seleção
de sensores e a exclusão exploratória de cinco coletas ocorreram após
examinar os mesmos dados usados na avaliação e podem elevar as estimativas
de desempenho. Os resultados devem ser tratados como prova de conceito
em condições controladas, sem estimativa validada de desempenho em novos
períodos, vasos independentes ou condições de campo.
